\section{Exponential products: a uniform excess estimate}\label{sec:exponential}
The exponential law supplies positive expected excess beyond every multiple
of the coefficient size. A Gaussian scale mixture makes that excess uniform
over all real coefficient lists and all shifts.
Take $W=Z$ to have the exponential distribution with density $e^{-z}$ on
$(0,\infty)$, and write $X_{\mathrm{exp}}=X_Z$.  Here $m_2=2$, so
Lemma~\ref{prob:first-moment} gives $\norm{y}\le3\E\sigma$.
Let $G$ denote a standard real Gaussian and put, for $K>0$,
\begin{equation}\label{prob:excess-constant}
 c(K)=\frac1{16}\E(|G|-2K)_+>0.
\end{equation}
The Laplace distribution as a Gaussian scale mixture belongs to the theory
of normal scale mixtures \cite{AndrewsMallows1974}; the characteristic-function
calculation below proves the exact identity needed here.
We prove the uniform conditional bound
\begin{equation}\label{prob:exponential-excess}
 \E\bigl((|y|-K\sigma)_+\mid\mathcal G_T,(B_i)_i\bigr)
 \ge c(K)\sigma.
\end{equation}
After conditioning as indicated, all coefficients, $d$, and $\sigma$ are
fixed.  Resample only the exponential multipliers to obtain $y'$ with the
same conditional law.  The triangle inequality and the positive-part
inequality give
\[
 2\E(|y|-K\sigma)_+
 \ge\E(|y-y'|-2K\sigma)_+.
\]
The difference of two independent standard exponentials has characteristic
function $(1+q^2)^{-1}$ at $q\in\R$.  This also equals
$\E\exp(-q^2U)$ for a standard exponential $U$, so that difference has the
law of $\sqrt{2U}\,G$, with $U$ and $G$ independent.  Using independent
such representations for the different multipliers, if $\sigma>0$ the
law of $(y-y')/\sigma$ is a centered Gaussian mixture with variance
\[
 V=\frac{2\sum_iB_i^2U_i}{\sigma^2},\qquad
 \E V=2,\qquad
 \E V^2=4\left(1+\frac{\sum_iB_i^4}{\sigma^4}\right)\le8.
\]
Since $\E[V\mathbf1_{\{V\ge1\}}]\ge1$, Cauchy--Schwarz gives
$\PP(V\ge1)\ge1/8$.  On that event the Gaussian excess above $2K$ is at
least $\E(|G|-2K)_+$.  Hence
\[
 \E(|y-y'|-2K\sigma)_+
 \ge\frac{\sigma}{8}\E(|G|-2K)_+,
\]
which proves~\eqref{prob:exponential-excess}.  For $\sigma=0$ the asserted
lower bound is zero, so the conclusion holds without normalization.


\begin{proposition}\label{prop:exponential}
For every $t,K>0$, the exponential product space satisfies
\begin{equation}\label{eq:exponential-parameter}
 \whdelta_{X_{\mathrm{exp}}}(t)
 \ge c(K)\left(\frac t3-\frac1K\right)_+.
\end{equation}
In particular, the exponential estimate in Theorem~\ref{thm:main} holds.
\end{proposition}
\begin{proof}
Fix $u\in E_T$ and a finite-product combination $y\in F_T$. The event
$A=\{|u|\le tK\sigma\}$ is fixed under conditioning on $\mathcal H_T$.
Since $\sigma\mathbf1_{A^c}\le |u|/(tK)$, the conditional excess estimate
gives
\begin{align*}
 M_t(u,y)
 &\ge t\E\bigl[\mathbf1_A(|y|-K\sigma)_+\bigr]
 \ge t c(K)\E(\sigma\mathbf1_A)\\
 &\ge t c(K)\E\sigma-\frac{c(K)}K\norm{u}
 \ge c(K)\left(\frac t3\norm{y}-\frac1K\norm{u}\right).
\end{align*}
Apply Lemma~\ref{lem:affine}. At $K=10/t$ this gives
$7t c(10/t)/30$, as asserted in Theorem~\ref{thm:main}.
\end{proof}

\begin{remark}[A fixed-error passage from finite tails]\label{rem:finite-head}
The excess estimate also permits a direct passage using a fixed positive
approximation tolerance, rather than letting the center error tend to zero
as in Lemma~\ref{lem:affine}.
Fix $t>0$, a finite predecessor-closed subtree $T$ containing the root,
$u\in E_T$, and a finite-product combination $y\in F_T$.
If $\norm{u}\le2$ and $\norm{y}\ge9/10$, then $\E\sigma\ge3/10$. On
$A=\{|u|\le10\sigma\}$, at least $1/10$ of the expected coefficient
size remains, because $\E(\sigma\mathbf1_{A^c})\le1/5$. Consequently
\eqref{prob:exponential-excess}, with $K=10/t$, gives
\[
 M_t(u,y)\ge \frac t{10}c(10/t).
\]
Finite-tail approximation extends this last bound to every $y\in F_T$
with $\norm{y}\ge1$. Put $L_t=t c(10/t)/10$. Approximating a unit center
with error less than $\min\{1,L_t/4\}$ ensures $\norm{u}<2$ and loses
less than $L_t/2$ by \eqref{prob:continuity}. Thus this passage gives the
weaker estimate $\whdelta_{X_{\mathrm{exp}}}(t)\ge t c(10/t)/20$ for every
$t>0$ directly, without using the affine conclusion of Lemma~\ref{lem:affine}.
\end{remark}
